\documentclass[11pt]{article}

\usepackage[margin=2.5cm]{geometry}
\usepackage{amsmath, amssymb}
\usepackage{mathtools}
\usepackage{xcolor}
\usepackage{framed}
\usepackage{hyperref}

\newcommand{\vect}[1]{\mathbf{#1}}

\title{\Large AI-solutions --- Lab 1: Modeling and Identification of the Parrot AR.Drone \\[2pt]
\large Questions 2.1, 2.2, 2.3}
\author{}
\date{}

\begin{document}

\maketitle

\begin{framed}
\textbf{\color{red} AI-SOLUTIONS -- NOT OFFICIAL.}
This document was generated by an AI assistant (Claude, Anthropic) as a self-check reference only.
It has not been reviewed by course staff and may contain sign or convention errors.
Use it exclusively to cross-check a solution worked out independently on paper, never as a submitted answer.
\end{framed}

\section*{Setup and assumptions}

We use the notation of the handout: $\vect{f}_T = \begin{bmatrix} f_1 & f_2 & f_3 & f_4 \end{bmatrix}^T$ for the individual rotor thrusts, $n_i = c f_i$ for the individual reaction torques, $l$ for the arm length (distance from the center of mass to each rotor), and $\vect{f}_g = mg\vect{e}_3$ for gravity expressed in $\{I\}$.

For Fig.~1 of the handout we take the following geometric convention, consistent with the figure:
\begin{itemize}
    \item Rotor 1 is located on the $+x_B$ axis, rotor 4 on the $+y_B$ axis, rotor 3 on the $-x_B$ axis, rotor 2 on the $-y_B$ axis (arms aligned with the body axes, as drawn in Fig.~1 of the handout).
    \item Every rotor thrust points along $-z_B$ (upward, since $z_B$ points down in the NED-aligned body frame).
    \item Rotors 1 and 3 spin one way and rotors 2 and 4 spin the opposite way (the two counter-rotating pairs mentioned in the statement), so that the reaction-torque contributions to yaw alternate in sign along the rotor index.
\end{itemize}
If your own numbering/sign convention differs, only the signs inside the matrices below change --- the derivation method is unaffected.

\section*{2.1 -- Expressions for $\vect{f}$ and $\vect{n}$}

\textbf{Translational force.} The total external force in $\{B\}$ is the sum of gravity (expressed in $\{I\}$, rotated into $\{B\}$) and the resultant of the four rotor thrusts (already aligned with $\{B\}$, all along $-z_B$):
\[
\vect{f} = R^T(\vect{\lambda})\,\vect{f}_g + N\vect{f}_T .
\]
Hence
\[
M = R^T(\vect{\lambda}) \in SO(3), \qquad
N = -\vect{e}_3\begin{bmatrix}1&1&1&1\end{bmatrix}
   = \begin{bmatrix} 0&0&0&0\\ 0&0&0&0\\ -1&-1&-1&-1 \end{bmatrix} \in \mathbb{R}^{3\times4}.
\]
($M$ depends on attitude because gravity is defined in $\{I\}$ and must be rotated into $\{B\}$; $N$ is constant because every thrust is already expressed in $\{B\}$ along the same axis.)

\textbf{Torque.} Each rotor $i$ produces (a) a moment from its thrust acting at moment arm $l$ about the two horizontal axes, and (b) a reaction torque $n_i = cf_i$ about $z_B$, alternating in sign with the spin direction. With the geometric convention above ($\vect{r}_i \times (-f_i\vect{e}_3)$ for each rotor position $\vect{r}_i$):
\[
n_x = l(f_2-f_4), \qquad n_y = l(f_1-f_3), \qquad n_z = c(f_1-f_2+f_3-f_4),
\]
so that
\[
\vect{n} = P\vect{f}_T, \qquad
P = \begin{bmatrix} 0 & l & 0 & -l \\ l & 0 & -l & 0 \\ c & -c & c & -c \end{bmatrix} \in \mathbb{R}^{3\times4}.
\]

\section*{2.2 -- Redefining $\{B\}$ with $x_B$ at $45^\circ$}

Rotating the body axes by $45^\circ$ about $z_B$ (new $x_B'$ midway between the arms of rotors 1 and 2) does \emph{not} change $\vect{f}$: every thrust still acts along $-z_B$ (unchanged, since the rotation is about $z_B$ itself), and $R^T(\vect{\lambda})$ is still just the attitude rotation expressed with respect to whichever axes are called $\{B\}$ -- its functional form does not depend on this relabeling. Hence
\[
M,\ N \ \text{unchanged}.
\]

$\vect{n} = P\vect{f}_T$ \emph{does} change, because the horizontal position of every rotor relative to the (now rotated) $x_B,y_B$ axes changes. The new $x_B'$ is midway between the old $+x_B$ (rotor 1's arm) and old $-y_B$ (rotor 2's arm), i.e.\ $x_B' = (x_B-y_B)/\sqrt2$, $y_B' = (x_B+y_B)/\sqrt2$. Expressing each rotor position in the new basis gives $\vect{r}_1' = \tfrac{l\sqrt2}{2}(1,1,0)$, $\vect{r}_2' = \tfrac{l\sqrt2}{2}(1,-1,0)$, $\vect{r}_3' = \tfrac{l\sqrt2}{2}(-1,-1,0)$, $\vect{r}_4' = \tfrac{l\sqrt2}{2}(-1,1,0)$ -- each rotor sits at $45^\circ$ from both axes, so each thrust now contributes to \emph{both} $n_x$ and $n_y$ (instead of purely one, as before). Using $\vect{r}\times(-f\vect{e}_3) = (-f r_y,\, f r_x,\, 0)$ for each rotor:
\[
n_x = \frac{l\sqrt2}{2}\big(f_2 + f_3 - f_1 - f_4\big), \qquad
n_y = \frac{l\sqrt2}{2}\big(f_1 + f_2 - f_3 - f_4\big),
\]
while $n_z$ is unchanged (yaw reaction torque depends only on spin direction, not on the axes' in-plane orientation):
\[
n_z = c(f_1-f_2+f_3-f_4).
\]
So
\[
P' = \begin{bmatrix}
-\frac{l\sqrt2}{2} & \frac{l\sqrt2}{2} & \frac{l\sqrt2}{2} & -\frac{l\sqrt2}{2} \\[4pt]
\frac{l\sqrt2}{2} & \frac{l\sqrt2}{2} & -\frac{l\sqrt2}{2} & -\frac{l\sqrt2}{2} \\[4pt]
c & -c & c & -c
\end{bmatrix} \neq P.
\]
This is exactly the standard ``$+$ configuration vs.\ $\times$ configuration'' distinction: in the $+$ configuration (2.1) each rotor loads a single rotational axis; in the $\times$ configuration (2.2, the real AR.Drone) every rotor loads both horizontal axes simultaneously, with a $\sqrt2/2$ geometric factor.

\section*{2.3 -- Input transformation $\vect{u} = L\vect{f}_T$}

Stacking $T=\sum_i f_i$ (from 2.1, first row all ones) on top of $\vect{n}=P\vect{f}_T$ (from 2.1) gives
\[
\vect{u} = \begin{bmatrix}T\\ \vect{n}\end{bmatrix} = L\vect{f}_T, \qquad
L = \begin{bmatrix} 1 & 1 & 1 & 1 \\ 0 & l & 0 & -l \\ l & 0 & -l & 0 \\ c & -c & c & -c \end{bmatrix}.
\]

\textbf{Usefulness.} $L$ is a constant, invertible matrix ($\det L \neq 0$ for $l,c\neq0$), so the map $\vect{f}_T \leftrightarrow \vect{u}$ is a lossless linear change of variables. Its point is that the rigid-body dynamics (1)--(3) only ever depend on $\vect{f}_T$ through exactly these four combinations: $T$ enters the translational dynamics (scaled along $z_B$) and $\vect{n}$ enters the rotational dynamics directly. The individual $f_i$ are physically coupled (each one affects both $T$ and one or two components of $\vect{n}$); $\vect{u}$ decouples the problem into four independent virtual channels -- one for altitude/thrust and three for the body-frame moments (roll, pitch, yaw). This lets the height and attitude controllers be designed directly in terms of $T$ and $\vect{n}$ (as done throughout the rest of the guide), with the physical per-rotor commands recovered only at the very last step via $L^{-1}$ -- i.e., $L$ is exactly the control-allocation / mixer stage of a real flight controller. The only caveat is that not every $\vect{u}$ is achievable, since $\vect{f}_T$ must remain nonnegative and bounded.

\textbf{Inverse transformation.} Solving $L\vect{f}_T=\vect{u}$ directly (summing/subtracting pairs of rows):
\[
f_1 = \frac{T}{4} + \frac{n_y}{2l} + \frac{n_z}{4c}, \qquad
f_2 = \frac{T}{4} + \frac{n_x}{2l} - \frac{n_z}{4c},
\]
\[
f_3 = \frac{T}{4} - \frac{n_y}{2l} + \frac{n_z}{4c}, \qquad
f_4 = \frac{T}{4} - \frac{n_x}{2l} - \frac{n_z}{4c},
\]
i.e. $\vect{f}_T = L^{-1}\vect{u}$ with
\[
L^{-1} = \begin{bmatrix}
1/4 & 0 & 1/(2l) & 1/(4c) \\
1/4 & 1/(2l) & 0 & -1/(4c) \\
1/4 & 0 & -1/(2l) & 1/(4c) \\
1/4 & -1/(2l) & 0 & -1/(4c)
\end{bmatrix}.
\]
Each rotor receives a quarter of the total thrust plus/minus the roll, pitch and yaw contributions needed to produce the requested $\vect{n}$ -- the familiar quadrotor mixer formula.

\end{document}
